\section*{Acknowledgements}
I wish to thank Peter Weisz for a critical reading
of the paper and for checking some of the formulae in sect.~2.
All numerical simulations were performed on a
dedicated PC cluster at CERN. I am grateful
to the CERN management for providing the required funds
and to the CERN IT Department for technical support.

\appendix

\section{Notational conventions}\label{app:A}

The Lie algebra $\sun$ of $\SUn$
may be identified with the linear space of all anti-hermitian
traceless $N\times N$ matrices.
With respect to a
basis $T^a$, $a=1,\ldots,N^2-1$, of such matrices,
the elements $X\in\sun$ are given by
$X=X^aT^a$ with real components $X^a$
(repeated group indices are automatically summed over).
The structure constants $f^{abc}$ in
the commutator relation
\begin{equation}
  [T^a,T^b]=f^{abc}T^c
  \label{eq:commutator}
\end{equation}
are real and totally anti-symmetric in the indices if
the normalization condition
\begin{equation}
  \tr\{T^aT^b\}=-\frac{1}{2}\delta^{ab}
  \label{eq:normalization}
\end{equation}
is imposed.
Moreover, $f^{acd}f^{bcd}=N\delta^{ab}$.

Gauge fields in the continuum theory
take values in the Lie algebra of the gauge group.
Their
normalization is usually such that
gauge transformations and covariant derivatives
do not involve the gauge coupling
(see, however, subsect.~2.2).
Lorentz indices $\mu,\nu,\ldots$ in $D$ dimensions
range from $0$ to $D-1$ and are automatically
summed over when they occur in matching pairs.
The space-time metric is assumed to be Euclidian. In particular,
$p^2=p_{\mu}p_{\mu}$ for any momentum $p$.

The lattice theories considered in this paper are
set up as usual on a four-dimen\-sional hyper-cubic
lattice with spacing $a$. In particular, lattice gauge fields $U(x,\mu)$
reside on the links
$(x,\mu)$ of the lattice and take values in the gauge group.
The link differential operators acting on functions
$f(U)$ of the gauge field are
\begin{equation}
  \partial^a_{x,\mu}f(U)=
  {\rmd\over\rmd s}f(\rme^{sX}U)\!\left.\vphantom{\frac{r}{s}}\right|_{s=0},
  \hskip1.7em
  X(y,\nu)=\begin{cases}T^a & \text{if $(y,\nu)=(x,\mu)$,}\\[1ex]
                        0  & \text{otherwise},\end{cases}
  \label{eq:linkderivatives}
\end{equation}
While these depend on the choice of the generators $T^a$,
the combination
\begin{equation}
  \partial_{x,\mu}f(U)=T^a\partial^a_{x,\mu}f(U)
  \label{eq:covlinkderivative}
\end{equation}
can be shown to be basis-independent.

\section{One-loop integrals}\label{app:B}

The Feynman integrals other than ${\cal E}_0$ which contribute to
$\langle E\rangle$ at order $g_0^4$ involve an integration over
two momenta, $q$ and $r$, and over at most two
flow-time parameters. In all cases the integrand is of the
form
\begin{equation}
  \rme^{-sq^2-ur^2-v(q+r)^2}{P(q,r)\over q^2r^2(q+r)^2}
\end{equation}
where $P(q,r)$ is a Lorentz-invariant polynomial and $s,u,v$ are linear
combinations of the flow-time parameters.

As already mentioned, these integrals can in principle be evaluated by
substituting the Feynman parameter representation for the propagators
$1/q^2$, $1/r^2$ and $1/(q+r)^2$.
A more economic computation is, however, always possible
using the basic integrals
\begin{align}
  \int_q\rme^{-sq^2}(q^2)^{-\alpha}&={s^{\alpha}\over(4\pi s)^{D/2}}
  {\Gamma(D/2-\alpha)\over\Gamma(D/2)},
  \\
  \int_{q,r}{\rme^{-sq^2-ur^2-v(q+r)^2}\over q^2}&=
  {2\over(4\pi)^D(D-2)(u+v)}(su+uv+vs)^{1-D/2}.
  \label{eq:basicintegralB}
\end{align}
The second formula, for example, allows the integral
\begin{equation}
  \int_{q,r}
  {\rme^{-t(q^2+r^2+(q+r)^2)}\over q^2r^2}=
  {1\over(4\pi)^4(2t)^2}\bigl\{8\ln2-4\ln3+\rmO(\eps)\bigr\}
\end{equation}
to be quickly evaluated once the propagator $1/r^2$ is replaced
by its Feynman parameter representation.

A special case are integrals like
\begin{equation}
\begin{split}
  \int_0^t\rmd s&\int_{q,r}
  {\rme^{-(t+s)(q^2+r^2)-(t-s)(q+r)^2}\over q^2r^2}\,
  qr\\
  &={1\over(4\pi)^4(2t)^2}\bigl\{\frac{1}{2}-4\ln2+2\ln3+\rmO(\eps)\bigr\},
\end{split}
\end{equation}
whose integrand is proportional to $qr$. Noting
\begin{equation}
  qr=\frac{1}{2}\left\{(q+r)^2-q^2-r^2\right\},
\end{equation}
the factor can be traded for a differentiation with respect to the
flow-time parameters. Most of the time, this allows one integral over
these parameters to be performed right away and thus leads to
simpler integrals without factors of $qr$.


\section{Numerical integration of the Wilson flow}\label{app:C}

On a finite lattice, the space ${\cal G}$ of all gauge fields
is a finite power of the gauge group and thus itself a Lie group.
The associated Lie algebra $\mathfrak{g}$ coincides with the linear
space of all link fields
with values in the Lie algebra of the gauge group. From this
abstract point of view,
the flow equation~\eqref{eq:wflow} is an ordinary first-order differential
equation of the form
\begin{equation}
   \dot{V_t}=Z(V_t)V_t,
   \label{eq:flowODE}
\end{equation}
where $V_t\in{\cal G}$ and $Z(V_t)\in\mathfrak{g}$.

The Runge--Kutta scheme
described in this appendix
obtains the solution of the flow equation at
times $t=n\eps$, $n=1,2,3,\ldots$, recursively,
starting from the initial configuration at $t=0$.
The rule for the integration
from time $t$ to $t+\eps$ is
\begin{align}
  W_0&=V_t,
  \nonumber\\
  W_1&=\exp\bigl\{\frac{1}{4}Z_0\bigr\}W_0,
  \nonumber\\
  W_2&=\exp\bigl\{\frac{8}{9}Z_1-\frac{17}{36}Z_0\bigr\}W_1,
  \nonumber\\
  V_{t+\eps}&=\exp\bigl\{\frac{3}{4}Z_2
                        -\frac{8}{9}Z_1+\frac{17}{36}Z_0\bigr\}W_2,
  \label{eq:RKstep}
\end{align}
where
\begin{equation}
  Z_i=\eps Z(W_i),\qquad i=0,1,2.
\end{equation}
Note that this rule is fully explicit.
Moreover, since the gauge field can be overwritten from one
equation to the next, and
since $Z_0$ can be overwritten by
$\frac{8}{9}Z_1-\frac{17}{36}Z_0$, intermediate
storage space for only one of these latter
fields is required.

A straightforward calculation
shows that the integration scheme~\eqref{eq:RKstep}
is accurate up to errors of order $\eps^4$ per step. The total
error of the integration up to a specified flow time thus scales
like $\eps^3$. Empirically one finds that
the integration is numerically stable in the
direction of positive flow time,
the integration errors in the link variables being
on the order of $10^{-6}$ if $\eps=0.01$.

\begin{thebibliography}{20}

\bibitem{Wilson}
K. G. Wilson, \textit{Confinement of quarks},
Phys. Rev. D10 (1974) 2445

\bibitem{ArnoldI}
V. I. Arnold,
\textit{Ordinary differential equations}, 3rd ed.\ (Springer-Verlag, Berlin, 2008)

\bibitem{Stout}
C. Morningstar, M. Peardon,
\textit{Analytic smearing of SU(3) link variables in lattice QCD},
Phys. Rev. D69 (2004) 054501

\bibitem{TrivMaps}
M. L\"uscher,
\textit{Trivializing maps, the Wilson flow and the HMC algorithm},
Commun. Math. Phys. 293 (2010) 899

\bibitem{BardeenEtAl}
W. A. Bardeen, A. J. Buras, D. W. Duke, T. Muta,
\textit{Deep inelastic scattering beyond the leading order
in asymptotically free gauge theories},
Phys. Rev. D18 (1978) 3998

\bibitem{SommerScale}
R. Sommer,
\textit{A new way to set the energy scale in lattice gauge theories
and its applications to the static force and $\alpha_s$ in
SU(2) Yang--Mills theory},
Nucl. Phys. B411 (1994) 839

\bibitem{GuagnelliEtAl}
M. Guagnelli, R. Sommer, H. Wittig (ALPHA collab.),
\textit{Precision computation of a low-energy reference scale in
quenched lattice QCD},
Nucl. Phys. B535 (1998) 389

\bibitem{DelDebbioTauQ}
L. Del Debbio, H. Panagopoulos, E. Vicari,
\textit{$\theta$-dependence of SU(N) gauge theories},
JHEP 0208 (2002) 044

\bibitem{StefanConference}
S. Schaefer, R. Sommer, F. Virotta,
\textit{Investigating the critical slowing down of QCD simulations},
PoS (LAT2009) 032

\bibitem{HairerEtAl}
E. Hairer, C. Lubich, G. Wanner,
\textit{Geometric numerical integration:
Structure-preserving algorithms for ordinary differential equations},
2nd ed.\ (Springer, Berlin, 2006)

\bibitem{LambdaParm}
S. Capitani, M. L\"uscher, R. Sommer, H. Wittig (ALPHA collab.),
\textit{Non-perturbative quark mass renormalization in quenched lattice QCD},
Nucl. Phys. B544 (1999) 669

\bibitem{RitbergenEtAl}
T. van Ritbergen, J. A. M. Vermaseren, S. A. Larin,
\textit{The four-loop $\beta$-function in Quantum Chromodynamics},
Phys. Lett. B400 (1997) 379

\bibitem{TopCharge}
M. L\"uscher,
\textit{Topology of lattice gauge fields},
Commun. Math. Phys. 85 (1982) 39

\bibitem{PhillipsStone}
A. Phillips, D. Stone, \textit{Lattice gauge fields,
principal bundles and the calculation of the topological charge},
Commun. Math. Phys. 103 (1986) 599

\bibitem{IndThm}
P. Hasenfratz, V. Laliena, F. Niedermayer,
\textit{The index theorem in QCD with a finite cutoff},
Phys. Lett. B427 (1998) 125

\bibitem{ChiGW}
L. Del Debbio, L. Giusti, C. Pica,
\textit{Topological susceptibility in SU(3) gauge theory},
Phys. Rev. Lett. 94 (2005) 032003

\bibitem{ChiInd}
M. L\"uscher,
\textit{Topological effects in QCD and the
problem of short distance singularities},
Phys. Lett. B593 (2004) 296

\bibitem{ChiProj}
L. Giusti, M. L\"uscher,
\textit{Chiral symmetry breaking and the Banks--Casher relation
in lattice QCD with Wilson quarks},
JHEP 0903 (2009) 013

\bibitem{ChiWI}
L. Giusti, G. C. Rossi, M. Testa, G. Veneziano,
\textit{The $U_A(1)$ problem on the lattice with Ginsparg--Wilson fermions},
Nucl. Phys. B628 (2002) 234

\bibitem{ChiWII}
L. Giusti, G. C. Rossi, M. Testa,
\textit{Topological susceptibility in full QCD with Ginsparg--Wilson fermions},
Phys. Lett. B587 (2004) 157

\end{thebibliography}



